% Part: sets-functions-relations
% Chapter: sets
% Section: basics

\documentclass[../../../include/open-logic-section]{subfiles}

\begin{document}

\olfileid{sfr}{set}{bas}
\olsection{Ekstensionalitas}

\emph{Himpunan} yaiku kumpulan objek kang dianggep {minangka} siji
objek. Objek-objek kang mbentuk himpunan kasebut diarani \emph{unsur}
utawa \emph{anggota} himpunan. Yen $x$ iku !!a{element} himpunan~$A$,
kita nulis $x \in A$; yen ora, kita nulis $x \notin A$.
Himpunan kang ora duwe !!{element}s diarani himpunan \emph{kosong}
lan dilambangake nganggo~"$\emptyset$".

\begin{explain}
Ora dadi prakara kepriye kita \emph{nemtokake} himpunan kasebut,
kepriye kita \emph{ngurutake} !!{element}sne, utawa malah \emph{kaping
pira} kita ngetung !!{element}sne. Kang wigati mung apa wae
!!{element}sne. Bab iki kita rumusake ing prinsip ing ngisor iki.
\end{explain}

\begin{defn}[Ekstensionalitas]
  Yen $A$ lan $B$ iku himpunan, mula $A = B$ yen lan mung yen
  saben !!{element} saka~$A$ uga !!a{element} saka~$B$, lan
  kosok baline.
\end{defn}

Ekstensionalitas mbenerake panganggone sawetara notasi. Umume, yen
kita duwe objek $a_{1}$, \dots, $a_{n}$, mula $\{a_{1}, \dots, a_{n}\}$
yaiku \emph{himpunan kasebut} kang !!{element}sne yaiku
$a_1, \ldots, a_n$. Kita negesake tembung "\emph{himpunan kasebut}",
amarga ekstensionalitas nerangake yen himpunan kaya mangkono mung
bisa ana \emph{siji}. Ekstensionalitas uga mbenerake pratelan iki:
  \[
    \{a, a, b\} = \{a, b\} = \{b,a\}.
  \]
Iki nuduhake yen, nalika ngrembug himpunan, kita ora nggatekake
urutane !!{element}s utawa kaping pira anggota kasebut disebutake.

\begin{tagblock}{novice}
\begin{ex}
Yen duwe sawenehing objek, objek kasebut bisa diklumpukake dadi
himpunan. Tuladhane, himpunan sedulure Richard ngemot siji wong,
lan bisa ditulis $S=\{\textrm{Ruth}\}$.
Himpunan wilangan wutuh positif kang kurang saka $4$ yaiku
$\{1, 2, 3\}$, nanging uga bisa ditulis $\{3, 2, 1\}$ utawa malah
$\{1, 2, 1, 2, 3\}$. Miturut ekstensionalitas, kabeh mau himpunan
kang padha. Amarga saben !!{element} saka $\{1, 2, 3\}$ uga
!!a{element} saka $\{3, 2, 1\}$ (lan saka $\{1,
2, 1, 2, 3\}$), lan kosok baline.
\end{ex}
\end{tagblock}

Asring kita nemtokake himpunan lumantar sawijining sipat kang
diduweni bebarengan dening !!{element}sne. Kanggo iku kita nganggo
notasi cekak iki: $\Setabs{x}{\phi(x)}$, kanthi $\phi(x)$
nglambangake sipat kang kudu diduweni~$x$ supaya bisa kalebu
ing antarane !!{element}s himpunan kasebut.

\begin{tagblock}{novice}
\begin{ex}
Ing tuladha kita, $S$ uga bisa ditemtokake minangka
\[
S = \Setabs{x}{x \text{ iku sedulure Richard}}.
\]
\end{ex}
\end{tagblock}

\begin{tagblock}{math}
\begin{ex}
Wilangan diarani \emph{sampurna} yen lan mung yen padha karo
gunggunge pembagi sejatine (yaiku, wilangan kang mbagi nganti entek
nanging ora padha karo wilangan kasebut). Tuladhane, $6$ iku
sampurna amarga pembagi sejatine yaiku $1$, $2$, lan~$3$, lan
$6 = 1 + 2 + 3$. Nyatane, $6$ iku siji-sijine wilangan wutuh
positif kang kurang saka $10$ lan sampurna. Mula, kanthi
ekstensionalitas, kita bisa nyatakake:
\[
	\{6\} = \Setabs{x}{x\text{ iku sampurna lan }0 \leq x \leq 10}
\]
Notasi ing sisih tengen diwaca "himpunan kabeh $x$ kang ndadekake
$x$ sampurna lan $0 \leq x \leq 10$". Pepadhan iki negesake yen,
nalika ngrembug himpunan, kita ora nggatekake carane himpunan kasebut
ditemtokake. Kanthi luwih umum, ekstensionalitas njamin yen tansah
mung ana siji himpunan kabeh $x$ kang ndadekake $\phi(x)$.
Mula, ekstensionalitas mbenerake anggone nyebut
$\Setabs{x}{\phi(x)}$ minangka \emph{himpunan kasebut} kang ngemot
kabeh $x$ kang ndadekake~$\phi(x)$.
\end{ex}
\end{tagblock}

Ekstensionalitas menehi cara kanggo nuduhake yen himpunan iku padha:
kanggo nuduhake $A = B$, tuduhna yen saben $x \in A$ uga $x \in B$,
lan saben $y \in B$ uga $y \in A$.

\begin{prob}
Buktekna yen himpunan kosong paling akeh ana siji, yaiku tuduhna
yen $A$ lan $B$ iku himpunan tanpa !!{element}s, mula $A = B$.
\end{prob}

\end{document}
